\section{Формалізована модель системи}

Для формалізації обрано стохастичну мережу Петрі з позначеннями за таблицею позначень~3.2
посібника~\cite{stetsenko-kursova}: позиція (кількість маркерів~$m$), дуга (кратність~$k$), перехід
(затримка~$d$, пріоритет~$r$, ймовірність запуску~$b$). Часткове захоплення спільної памʼяті
задається кратністю дуг, конкуренція подій за ресурс -- спільною вхідною позицією, а стратегія
упорядкування -- пріоритетами та ймовірностями запуску. Мережу побудовано й перевірено в
редакторі й симуляторі PetriObjModelPaint~\cite{petriobjmodelpaint}.

\subsection{Стохастична мережа Петрі моделі}

Подіям концептуальної моделі (рисунок~\ref{fig:process}) відповідають переходи: надходження
й виконання завдання типу~$i$ і резервування памʼяті для завдання~C (рисунок~\ref{fig:petri};
штрихова стрілка -- інформаційна дуга~\cite{stetsenko-modeling}, яка перевіряє наявність
маркера, але не вилучає його;
кратність дуг без числа дорівнює~1). Позиції наведено в таблиці~\ref{tab:places}, параметри
переходів -- у таблиці~\ref{tab:transitions} ($i\in\{A,B,C\}$ -- тип завдання).

\begin{figure}
  \centering
  % Рисунок 2.1 — стохастична мережа Петрі моделі B14 (розкладка відповідає model/B14.pnml)
\begin{tikzpicture}[
    font=\setstretch{1}\fontsize{11}{13}\selectfont, >={Stealth[length=2mm]},
    pl/.style={draw, circle, minimum size=8mm, inner sep=0pt},
    tr/.style={draw, fill=black!12, rectangle, minimum width=2.2mm, minimum height=9mm, inner sep=0pt},
    arc/.style={->},
    rd/.style={->, dashed},
    nm/.style={font=\setstretch{1}\fontsize{10}{12}\selectfont},
    par/.style={font=\setstretch{1}\fontsize{10}{12}\selectfont, align=center},
    w/.style={font=\setstretch{1}\fontsize{10}{12}\selectfont, fill=white, inner sep=1pt},
  ]
  \def\yA{0} \def\yB{-2.6} \def\yC{-6.0}
  \def\xS{7.2}
  % Джерела та підчерги
  \foreach \i/\y/\l in {A/\yA/\lambda_A, B/\yB/\lambda_B, C/\yC/\lambda_C}{
    \node[pl, label={[nm]above:Gen\i}] (Gen\i) at (0,\y) {1};
    \node[tr, label={[nm]above:Arr\i}, label={[par]below:$\exp(1/\l)$}] (Arr\i) at (1.8,\y) {};
    \node[pl, label={[nm]above:Q\i}] (Q\i) at (3.6,\y) {};
    \draw[arc] ([yshift=1.2mm]Gen\i.east) -- ([yshift=1.2mm]Arr\i.west);
    \draw[arc] ([yshift=-1.2mm]Arr\i.west) -- ([yshift=-1.2mm]Gen\i.east);
    \draw[arc] (Arr\i) -- (Q\i);
  }
  % Виконання
  \node[tr, label={[nm]above:ServA}, label={[par]below left:$\exp(4)$; $r_A$, $b_A$}] (ServA) at (\xS,\yA) {};
  \node[tr, label={[nm]above right:ServB}, label={[par]below left:$\exp(5)$; $r_B$, $b_B$}] (ServB) at (\xS,\yB) {};
  \node[tr, label={[nm]below left:ServC}, label={[par]below right:$\exp(10)$}] (ServC) at (\xS,\yC) {};
  % Виходи
  \node[pl, label={[nm]above:OutA}] (OutA) at (14.3,\yA) {};
  \node[pl, label={[nm]left:OutB}] (OutB) at (5.8,-1.3) {};
  \node[pl, label={[nm]above:OutC}] (OutC) at (14.3,\yC) {};
  % Памʼять і резервування
  \node[pl, label={[nm]right:Mem}] (Mem) at (12.5,\yB) {2};
  \node[tr, label={[nm]above right:ResC}, label={[par]above left:$0$; $r_C$}] (ResC) at (5.0,-4.4) {};
  \node[pl, label={[nm]right:RC}] (RC) at (6.1,-5.25) {};
  % Дуги
  \draw[arc] (QA) -- (ServA);
  \draw[arc] (QB) -- (ServB);
  \draw[arc] (QC) -- (ServC);
  \draw[arc] (ServA) -- (OutA);
  \draw[arc] (ServB) -- (OutB);
  \draw[arc] (ServC) -- (OutC);
  \draw[arc] (Mem) to[bend left=5] (ServA);
  \draw[arc] (ServA) to[bend left=5] (Mem);
  \draw[arc] ([yshift=1.2mm]ServB.east) -- ([yshift=1.2mm]Mem.west);
  \draw[arc] ([yshift=-1.2mm]Mem.west) -- ([yshift=-1.2mm]ServB.east);
  \draw[arc] (Mem) -- node[w, pos=0.45, below] {$k_C$} (ResC);
  \draw[arc] (ResC) -- node[w, above right] {$k_C$} (RC);
  \draw[arc] (RC) -- node[w, below left] {2} (ServC);
  \draw[arc] (ServC) -- node[w, pos=0.5, below right] {2} (Mem);
  \draw[rd] (QC) -- (ResC);
\end{tikzpicture}

  \caption{Стохастична мережа Петрі інформаційної системи обробки завдань}
  \label{fig:petri}
\end{figure}

\begin{kpitable}{|c|L{12.4cm}|c|}{Позиції мережі Петрі та їх початкове маркірування~$m$\label{tab:places}}
  {\textbf{Позиція} & \textbf{Зміст} & \textbf{$m$}}
  Gen$i$ & умова роботи джерела типу~$i$ (маркер повертає перехід Arr$i$) & 1 \\ \hline
  Q$i$ & підчерга завдань типу~$i$ (кількість маркерів -- довжина підчерги) & 0 \\ \hline
  Mem & вільні одиниці оперативної памʼяті & $M=2$ \\ \hline
  RC & одиниці памʼяті, зарезервовані для завдання~C, що очікує & 0 \\ \hline
  Out$i$ & виконані завдання типу~$i$ (лічильник виходів із системи) & 0 \\ \hline
\end{kpitable}

\begin{kpitable}{|c|L{8.6cm}|L{2.6cm}|c|c|}{Переходи мережі Петрі\label{tab:transitions}}
  {\textbf{Перехід} & \textbf{Подія} & \textbf{$d$, хв} & \textbf{$r$} & \textbf{$b$}}
  ArrA & надходження завдання A & $\exp(5/\gamma)$ & 0 & 1 \\ \hline
  ArrB & надходження завдання B & $\exp(14{,}29/\gamma)$ & 0 & 1 \\ \hline
  ArrC & надходження завдання C & $\exp(16{,}67/\gamma)$ & 0 & 1 \\ \hline
  ServA & виконання завдання A & $\exp(4)$ & $r_A$ & $b_A$ \\ \hline
  ServB & виконання завдання B & $\exp(5)$ & $r_B$ & $b_B$ \\ \hline
  ResC & резервування памʼяті для завдання~C & 0 & $r_C$ & 1 \\ \hline
  ServC & виконання завдання C & $\exp(10)$ & 0 & 1 \\ \hline
\end{kpitable}

Затримки Arr$i$ задано через~$\gamma$ ($\lambda_i=\gamma\lambda_i^0$), Serv$i$ -- середніми $1/\mu_i$
з умови; $r_A$, $r_B$, $r_C$, $b_A$, $b_B$ і~$k_C$ залежать від стратегії (підрозділ~\ref{sec:strategies-net}).

Затримку $\exp(a)$ генерують методом оберненої функції~\cite{stetsenko-modeling}:
\[
  d = -a\ln(1-u), \qquad u \sim U[0;\,1),
\]
де $u$ -- псевдовипадкове число, $a = 1/\lambda_i$ для Arr$i$ і $1/\mu_i$ для Serv$i$. Якщо готові
кілька переходів з однаковим найвищим пріоритетом, генерують $u \sim U[0;\,1)$ і запускають перший
перехід $T_j$ групи, для якого $u < b_1 + \cdots + b_j$ (перехід з $b=1$ у групі з $n$ переходів
отримує ймовірність $1/n$).

Перехід Arr$i$ додає маркер у Q$i$. Перехід ServA (ServB) запускається за наявності маркерів у
QA (QB) і в Mem, утримує їх протягом виконання, повертає одиницю памʼяті в Mem і
передає завдання в OutA (OutB). Переходи виконання багатоканальні: за двох вільних одиниць
одночасно виконуються два завдання A/B (пари AA, AB, BB). ServC стартує лише тоді, коли для C
зарезервовано обидві одиниці (кратність дуги з~RC -- 2), і повертає їх у Mem після виконання;
поки виконується C, жодне інше завдання стартувати не може -- C монополізує систему.

Кількість завдань типу~$i$, що виконуються, дорівнює кількості активних каналів Serv$i$, і в
будь-який момент справджується інваріант памʼяті
\begin{equation}\label{eq:mem-invariant}
  m_{\text{Mem}} + m_{\text{RC}} + k_C\,n_{\text{ResC}} + n_{\text{ServA}} + n_{\text{ServB}}
    + 2\,n_{\text{ServC}} = M = 2,
\end{equation}
де $m_P$ -- кількість маркерів у позиції~$P$, $n_T$ -- кількість активних каналів переходу~$T$
(для ResC ненульова лише миттєво, бо $d=0$). Інваріант гарантує відсутність станів, заборонених
умовою, і використовується у верифікації (розділ~3).

\subsection{Відтворення стратегій упорядкування черги}\label{sec:strategies-net}

Пріоритет переходу визначає лише, яка з одночасно готових подій відбудеться першою. Цього
досить, щоб завдання вищого пріоритету першим займало памʼять, що звільнилася, але недостатньо
для правила резервування (підрозділ~\ref{sec:process}): коли завдання~C вищого пріоритету чекає двох вільних
одиниць памʼяті, а вільна лише одна, переходи ServA і ServB залишаються готовими і займали б її.
Тому до мережі додано перехід ResC і позицію RC. Перехід ResC готовий, поки в підчерзі QC є
завдання (інформаційна дуга) і в Mem є щонайменше $k_C$ вільних одиниць; запускаючись, він
переносить $k_C$ одиниць памʼяті в RC, де вони недоступні для A і~B. Отже, правило
резервування відтворено стандартними засобами формалізму: пріоритетом, кратністю дуг та
інформаційною дугою.

Параметри стратегій наведено в таблиці~\ref{tab:strategy-params} (більше~$r$ -- вищий пріоритет).

\needspace{10\baselineskip}
\begin{kpitable}{|c|c|c|c|c|c|c|}{Параметри мережі для стратегій упорядкування\label{tab:strategy-params}}
  {\textbf{$s$} & \textbf{Упорядкування} & \textbf{$r_A$} & \textbf{$r_B$} & \textbf{$r_C$} &
   \textbf{$k_C$} & \textbf{$b_A$; $b_B$}}
  1 & A $>$ B $>$ C   & 3 & 2 & 1 & 2 & 1; 1 \\ \hline
  2 & B $>$ A $>$ C   & 2 & 3 & 1 & 2 & 1; 1 \\ \hline
  3 & C $>$ B $>$ A   & 1 & 2 & 3 & 1 & 1; 1 \\ \hline
  4 & C $>$ A $>$ B   & 2 & 1 & 3 & 1 & 1; 1 \\ \hline
  5 & C $>$ (A $=$ B) & 1 & 1 & 2 & 1 & $0{,}5$; $0{,}5$ \\ \hline
\end{kpitable}

У стратегіях 3--5 перехід ResC має найвищий пріоритет і кратність $k_C=1$. Поки завдання~C
очікує, кожна одиниця памʼяті, що звільнилася, одразу резервується для нього, тому A і~B не
стартують навіть за наявності вільної одиниці; щойно в RC накопичуються дві одиниці, стартує
ServC. За інваріантом~\eqref{eq:mem-invariant} у RC не може бути більше двох маркерів. Якщо
завдання~C надходить за вільної памʼяті, ResC спрацьовує двічі в той самий момент часу і C
стартує без очікування.

У стратегіях 1--2 перехід ResC має найнижчий пріоритет і кратність $k_C=2$, тобто резервує
памʼять лише тоді, коли обидві одиниці вільні. Якби в цей момент у QA чи QB було завдання,
раніше спрацював би перехід ServA чи ServB з вищим пріоритетом, тому C стартує лише за порожніх
підчерг A і~B, як вимагає правило. Кратність $k_C=2$ тут обовʼязкова: за $k_C=1$ завдання~C
утримувало б одну одиницю, поки виконується інше завдання, і блокувало б A і~B, які
надходять пізніше, -- тобто отримало б фактичний пріоритет над ними.

Порядок A і~B задають $r_A$ і $r_B$; у стратегії~5 вони рівні, і конфлікт ServA і ServB за
вільну одиницю памʼяті розвʼязують ймовірності $b_A=b_B=0{,}5$; якщо готовий лише один із цих
переходів, він спрацьовує без розіграшу. Обслуговування без переривання
забезпечене тим, що запущений перехід утримує маркери до кінця затримки.

\subsection{Обчислення вихідних характеристик}

Статистику збирають на інтервалі $[T_{\text{п}};\,T_{\text{мод}}]$ тривалістю
$\tau=T_{\text{мод}}-T_{\text{п}}$. Якщо $t_j$ -- моменти подій, $\Delta t_j=t_{j+1}-t_j$, а
$m_P(t_j)$ і $n_T(t_j)$ -- стан протягом $j$-го інтервалу, то середні за часом
\begin{equation}\label{eq:time-means}
  \bar m_P = \frac{1}{\tau}\sum_j m_P(t_j)\,\Delta t_j, \qquad
  \bar n_T = \frac{1}{\tau}\sum_j n_T(t_j)\,\Delta t_j .
\end{equation}

Середня кількість завдань типу~$i$ у черзі та в системі, фактична інтенсивність надходження
($N_{\text{Arr}i}$ -- кількість спрацьовувань Arr$i$ за час~$\tau$) і, за законом
Літтла~\cite{kleinrock}, що для стаціонарного режиму не залежить від порядку обслуговування,
середні часи очікування й перебування:
\begin{equation}\label{eq:L-lambda}
  L_{q,i} = \bar m_{\text{Q}i}, \qquad
  L_i = \bar m_{\text{Q}i} + \bar n_{\text{Serv}i}, \qquad
  \hat\lambda_i = \frac{N_{\text{Arr}i}}{\tau}.
\end{equation}
\begin{equation}\label{eq:W-type}
  W_{q,i} = \frac{L_{q,i}}{\hat\lambda_i}, \qquad W_i = \frac{L_i}{\hat\lambda_i}.
\end{equation}
Загальні показники -- середні за всіма завданнями, тобто зважені за кількістю завдань кожного
типу (для $W$ маємо $\sum_i \hat\lambda_i W_i / \sum_i \hat\lambda_i = \sum_i L_i/\sum_i\hat\lambda_i$),
а допоміжний коефіцієнт використання памʼяті~\eqref{eq:U} -- частка зайнятих одиниць без
зарезервованих:
\begin{equation}\label{eq:W-total}
  L_q = \sum_{i} L_{q,i}, \qquad
  W_q = \frac{\sum_i L_{q,i}}{\sum_i \hat\lambda_i}, \qquad
  W = \frac{\sum_i L_i}{\sum_i \hat\lambda_i}.
\end{equation}
\begin{equation}\label{eq:U}
  U = \frac{\bar n_{\text{ServA}} + \bar n_{\text{ServB}} + 2\,\bar n_{\text{ServC}}}{M}
    = 1 - \frac{\bar m_{\text{Mem}} + \bar m_{\text{RC}}}{M}.
\end{equation}
Друга рівність випливає з інваріанту~\eqref{eq:mem-invariant}; $\bar m_{\text{RC}}/M$ -- частка
памʼяті, що простоює через резервування. Формули~\eqref{eq:time-means}--\eqref{eq:W-total} дають
саме ті характеристики, які названо в умові задачі: середній час перебування завдання в системі
($W$, $W_i$), середню кількість завдань у черзі ($L_q$, $L_{q,i}$) і середній час очікування в
черзі ($W_q$, $W_{q,i}$). Вони мають сенс лише за стійкого режиму, коли довжини черг не
зростають необмежено; кожну оцінюють як середнє
за $p$ прогонами з точністю 2~\% середнього і довірчою ймовірністю 0,95 (підрозділ~\ref{sec:exp-runs}).

\subsection{Висновки до розділу}

Систему формалізовано стохастичною мережею Петрі з 11 позицій і 7 переходів: обмеження на
одночасне виконання відтворено позицією памʼяті й кратністю дуг, стратегії -- пріоритетами,
ймовірностями запуску й кратністю дуги резервування, правило резервування -- переходом ResC з
інформаційною дугою. Наведено генератори, формули відгуків і інваріант памʼяті.
